\section{Metric Alignment and Render Supervision}
\label{sec:render_supervision}

The data generator is the main methodological contribution.  Its purpose is
not to create arbitrary synthetic courts; it turns a reconstructed real tennis
complex into a coordinate-aware renderer whose images and labels share one
auditable geometric source of truth.

\subsection{Frames and public reconstruction boundary}
\label{sec:frames}

We write \(\T{b}{a}\) for a homogeneous transform that maps a point expressed
in frame \(a\) into frame \(b\).  The public reconstruction export uses a
right-handed normalized scene frame \(\Scene_n\).  Cameras use the OpenCV axes
\(+x\) right, \(+y\) down, and \(+z\) forward, and publish
\(\T{\Scene_n}{\Cam}\).  A metric adapter defines a reciprocal uniform scale
between \(\Scene_n\) and a metre-valued scene frame \(\Scene_m\).  Each
regulation court \(i\) has a right-handed frame \(\Court_i\), in metres, with
\(+X\) toward the right sideline, \(+Y\) toward the far baseline, and \(+Z\)
up.  These conventions are summarized in \cref{tab:coordinate_contracts}.

The reconstruction program owns structure from motion, Gaussian training, and
rasterization.  We communicate with it only through two public commands: one
exports a validated standard scene, and the other renders a supplied camera.
The export validator checks schemas, camera identifiers, shapes, dtypes,
finite values, intrinsics, proper rotations, and reciprocal transforms before
alignment begins.  This boundary prevents the dataset from depending on an
unstated checkpoint or coordinate convention.

\begin{table}[t]
  \centering
  \caption{Coordinate frames.  All transforms are stored as proper homogeneous
  matrices; no quaternion convention is required by the data generator.}
  \label{tab:coordinate_contracts}
  \small
  \begin{tabularx}{\linewidth}{@{}p{0.15\linewidth}p{0.20\linewidth}X@{}}
    \toprule
    Frame & Units / axes & Contract \\
    \midrule
    NHT scene \(\Scene_n\) & Normalized, right-handed & Public reconstruction
      and render coordinates; Gaussian means and captured cameras live here. \\
    Metric scene \(\Scene_m\) & Metres, right-handed & Reciprocal uniform
      scaling of \(\Scene_n\); common carrier for all accepted courts. \\
    Court \(\Court_i\) & Metres; \(+X\) right sideline, \(+Y\) far baseline,
      \(+Z\) up & Regulation geometry and physical point indices. \\
    Camera \(\Cam\) & OpenCV; \(+x\) right, \(+y\) down, \(+z\) forward &
      \(\T{\Scene}{\Cam}\) is camera-to-scene; projection uses its inverse. \\
    Canonical court \(\Canon\) & Metres, target-relative & Court frame after
      camera-end choice \(S\in\{I,R_z(\pi)\}\); pose10D authority. \\
    Model image & Augmented pixels & KP logits, principal point, focal length,
      and reprojection residual are all expressed in this same frame. \\
    \bottomrule
  \end{tabularx}
\end{table}



\subsection{From image lines to a metric court layout}
\label{sec:alignment}

\begin{figure}[t]
  \centering
  \resizebox{\linewidth}{!}{% メートル系アライメントの流れ。すべて本稿のために作成した TikZ 図である。
\begin{tikzpicture}[
  font=\sffamily\scriptsize,
  frame/.style={draw=black!28, rounded corners=2.5pt, fill=none,
    inner sep=3mm, line width=0.45pt},
  title/.style={font=\sffamily\bfseries\footnotesize, text=black!85},
  tinylabel/.style={font=\sffamily\fontsize{6.3}{7.2}\selectfont,
    text=black!70, align=center},
  fitcam/.style={regular polygon, regular polygon sides=3, minimum size=3.2mm,
    inner sep=0pt, draw=gsblue, fill=gsblue!20, line width=0.5pt},
  holdcam/.style={regular polygon, regular polygon sides=3, minimum size=3.2mm,
    inner sep=0pt, draw=poseorange, fill=poseorange!22, line width=0.5pt},
  arr/.style={-{Stealth[length=1.7mm,width=1.2mm]}, line width=0.65pt},
  dot/.style={circle, inner sep=0pt, minimum size=1.2mm},
  gate/.style={rounded corners=1.5pt, inner sep=2pt, draw=lossred,
    fill=lossred!8, font=\sffamily\fontsize{6.2}{7.2}\selectfont, align=center},
]

  % ---------- (a) カメラ証拠を fit / holdout に先に分ける ----------
  \begin{scope}[local bounding box=pa]
    \node[fitcam, shape border rotate=-90] (f1) at (0.05,2.05) {};
    \node[fitcam, shape border rotate=-90] (f2) at (1.05,2.25) {};
    \node[fitcam, shape border rotate=-90] (f3) at (2.05,2.05) {};
    \node[holdcam, shape border rotate=-90] (h1) at (1.05,0.05) {};
    \draw[fill=softgray!70, draw=black!30, line width=0.4pt]
      (-0.35,0.60) rectangle (2.45,1.55);
    \node[tinylabel] at (1.05,1.08) {コート面の画像領域};
    \draw[arr, draw=gsblue] (f1) -- (0.40,1.55);
    \draw[arr, draw=gsblue] (f2) -- (1.05,1.55);
    \draw[arr, draw=gsblue] (f3) -- (1.70,1.55);
    \draw[arr, draw=poseorange, dashed] (h1) -- (1.05,0.60);
    \node[tinylabel, text=gsblue, anchor=south] at (1.05,2.50)
      {fit カメラ（変換の推定）};
    \node[tinylabel, text=poseorange, anchor=north] at (1.05,-0.25)
      {holdout カメラ（凍結検証）};
  \end{scope}

  % ---------- (b) 地面推定と光線交差 ----------
  \begin{scope}[shift={(4.6,0)}, local bounding box=pb]
    \coordinate (bl) at (0,-0.30);
    \coordinate (br) at (3.3,-0.30);
    \coordinate (tl) at (0.55,1.05);
    \coordinate (tr) at (2.75,1.05);
    \fill[softgray!80] (bl) -- (br) -- (tr) -- (tl) -- cycle;
    \draw[black!35, line width=0.4pt] (bl) -- (br) -- (tr) -- (tl) -- cycle;
    \draw[black!18] (0.50,-0.30) -- (0.95,1.05)
      (1.15,-0.30) -- (1.45,1.05)
      (1.85,-0.30) -- (1.90,1.05)
      (2.55,-0.30) -- (2.35,1.05);
    \node[fitcam, shape border rotate=-90] (g1) at (0.15,2.20) {};
    \coordinate (e1) at (1.05,0.30);
    \coordinate (e2) at (1.95,-0.20);
    \draw[arr, draw=gsblue] (g1) -- (e1);
    \draw[arr, draw=gsblue, dashed] (g1) -- (e2);
    \node[dot, fill=gsblue] at (e1) {};
    \node[dot, fill=gsblue] at (e2) {};
    \node[tinylabel, text=gsblue, anchor=south west] at (0.45,2.35)
      {線画素の光線と\\推定地上面の交差};
    \node[tinylabel, anchor=north, text width=42mm] at (1.65,-0.55)
      {採用画素数と光線--平面角が\\閾値未満なら、そのカメラを外す};
  \end{scope}

  % ---------- (c) 共通尺度の 2 面当てはめ ----------
  \begin{scope}[shift={(9.7,0)}, local bounding box=pc]
    \draw[courtgreen, line width=0.6pt, fill=courtgreen!7]
      (0.20,0.85) rectangle (1.75,2.05);
    \draw[courtgreen, line width=0.35pt] (0.975,0.85) -- (0.975,2.05);
    \draw[black!28, line width=0.35pt] (0.20,1.45) -- (1.75,1.45);
    \draw[courtgreen, line width=0.6pt, fill=courtgreen!7]
      (2.50,0.55) rectangle (4.05,1.75);
    \draw[courtgreen, line width=0.35pt] (3.275,0.55) -- (3.275,1.75);
    \draw[black!28, line width=0.35pt] (2.50,1.15) -- (4.05,1.15);
    \draw[<->, courtgreen, line width=0.5pt] (0.35,2.20) -- (3.90,2.20);
    \node[tinylabel, text=courtgreen, anchor=south] at (2.125,2.22)
      {共通尺度 \(s\)（メートル換算）};
    \node[gate, anchor=north] at (2.125,0.40)
      {中心間距離 \(\geq 10.97\) m、\\フットプリント重なり \(\leq 10^{-9}\)};
    \node[tinylabel, anchor=north, text width=44mm] at (2.125,-0.55)
      {トポロジーと尺度が食い違えば\\シーン全体を棄却する};
  \end{scope}

  % ---------- パネル枠と見出し ----------
  \node[frame, fit=(pa)] (fa) {};
  \node[frame, fit=(pb)] (fb) {};
  \node[frame, fit=(pc)] (fc) {};
  \node[title, anchor=south] at ($(fa.north)+(0,1.3mm)$) {(a) 証拠の分割};
  \node[title, anchor=south] at ($(fb.north)+(0,1.3mm)$) {(b) 地上面への持ち上げ};
  \node[title, anchor=south] at ($(fc.north)+(0,1.3mm)$) {(c) 共通尺度の 2 面当てはめ};

  \draw[arr] (fa.east) -- (fb.west);
  \draw[arr] (fb.east) -- (fc.west);
\end{tikzpicture}
}
  \caption{\textbf{Fail-closed 3DGS-to-court alignment.}  Captured cameras are
  partitioned before fitting.  Court-line pixels from fit cameras are intersected
  with a robust ground plane.  Two regulation templates are fitted with one
  shared scene-units-per-metre scale, followed by rigid fitting in metric scene
  coordinates.  The frozen transform must explain held-out cameras and full
  court topology.  Camera-end canonicalization uses only the two proper
  rotations \(I\) and \(R_z(\pi)\).}
  \label{fig:alignment}
\end{figure}

\paragraph{Ground plane and lifted evidence.}
Sparse scene points provide candidate ground support.  We estimate
\(
  \pi_g=\{\mathbf{x}:\mathbf{n}^{\top}\mathbf{x}+b=0\}
\)
with RANSAC, refine its support three times, and reject planes whose normal,
support, bounds, or camera-side tests are implausible.  For a detected line
pixel \(\tilde{\mathbf{u}}\), its captured camera defines the scene ray
\(
  \mathbf{r}(\lambda)=\mathbf{o}+\lambda\mathbf{d}
\).
The intersection
\begin{equation}
  \lambda^*=-\frac{\mathbf{n}^{\top}\mathbf{o}+b}
                    {\mathbf{n}^{\top}\mathbf{d}}
  \label{eq:ray_plane}
\end{equation}
is retained only when the ray--plane angle, positive distance, and complex
bounds pass strict gates.  Cameras with no observable line evidence are
excluded explicitly; insufficient remaining evidence is an error rather than
a fallback.

\paragraph{Shared-scale template hypotheses.}
The lifted points form a measured evidence set \(\mathcal{E}\) in a 2D basis
on the ground plane.  We sample every regulation line segment of a canonical
metre-valued template \(\mathcal{Q}\).  With
\(d(\mathbf{a},\mathcal{E})=\min_{\mathbf{e}\in\mathcal{E}}
\|\mathbf{a}-\mathbf{e}\|_2\), the implemented proposal score is
\begin{equation}
  J(\mathbf{c},\theta,s)=\frac{1}{|\mathcal{Q}|}
  \sum_{\mathbf{q}\in\mathcal{Q}}
  \exp\!\left[-\frac{1}{2}
    \left(\frac{d(s\mathbf{R}(\theta)\mathbf{q}+\mathbf{c},\mathcal{E})}
                    {s\sigma}\right)^2\right],
  \qquad
  (\mathbf{c}^*,\theta^*,s^*)=\arg\max J,
  \label{eq:template_fit}
\end{equation}
where \(s\) is measured in NHT scene units per metre and the configured score
width is \(\sigma=0.5\) m.  A bounded differential-evolution search covers tiled
centres, orientation bands, and scale.  After fitting the first court, its
evidence is suppressed before the second hypothesis is searched.  The two
hypotheses must have compatible line topology, negligible footprint overlap,
adequate centre separation, and native scales that agree within the declared
tolerance.  We then refit their centres at one common scale.  This multi-court
constraint resolves metric scale before any rigid court transform is fitted.

Let \(s_*\) be the accepted common scale.  The metric adapter is
\begin{equation}
  \mathbf{x}_{\Scene_n}=s_*\mathbf{x}_{\Scene_m},\qquad
  \mathbf{x}_{\Scene_m}=s_*^{-1}\mathbf{x}_{\Scene_n}.
  \label{eq:metric_adapter}
\end{equation}
For matched canonical court points \(\mathbf{x}_j\) and lifted metric-scene
points \(\mathbf{y}_j\), the final transform is the Kabsch solution
\begin{equation}
  (\mathbf{R}_i,\mathbf{t}_i)=
  \arg\min_{\mathbf{R}\in\SOthree,\mathbf{t}\in\R^3}
    \sum_j \left\|\mathbf{R}\mathbf{x}_j+\mathbf{t}-\mathbf{y}_j\right\|_2^2,
  \quad
  \T{\Scene_m}{\Court_i}=
  \begin{bmatrix}\mathbf{R}_i&\mathbf{t}_i\\\mathbf{0}^{\top}&1\end{bmatrix}.
  \label{eq:kabsch}
\end{equation}
Reflection is corrected so that \(\det\mathbf{R}_i=1\).  Scale is absent from
\cref{eq:kabsch} because it was already resolved by \cref{eq:metric_adapter}.

\paragraph{Disjoint verification.}
Camera ownership is fixed before fitting.  Correspondences from fit cameras
estimate \cref{eq:kabsch}; holdout cameras evaluate that frozen transform
without refitting.  In addition to nearest matched points, the verifier samples
the entire court template so that a good residual on one visible baseline
cannot hide a globally incorrect court.  Per-segment coverage, two-court
topology, scale agreement, and reciprocal-transform checks must all pass.  The
selected production thresholds in \cref{tab:alignment_gates} are executable
acceptance conditions, not reported benchmark accuracy; the complete typed
configuration remains the authority.

\begin{table}[t]
  \centering
  \caption{\textbf{採用したアライメント判定の閾値。}
  これらは実行可能な受理条件であり、
  達成された精度ではない。実際の権威は型付き設定ファイル側にあり、
  本表はその要約である。
  完全な一覧は付録~\ref{app:thresholds} に置く。
  }
  \label{tab:alignment_gates}
  \small
  \begin{tabularx}{\linewidth}{@{}p{0.32\linewidth}X@{}}
    \toprule
    判定 & 閾値 \\
    \midrule
    証拠の所有権 & 決定的に選ぶ 48 カメラ。
      fit 8 : holdout 4 の単位で固定分割（48 カメラでは 32 / 16） \\
    地上面 & 支持点 \(\geq 1000\)、
      normal と上方向の余弦 \(\geq 0.97\)、
      正のカメラ側割合 \(=1\) \\
    投影証拠 & 有効なカメラあたり \(\geq 20\) 点 \\
    コート仮説 & ちょうど 2 面。中心間距離 \(\geq 10.97\) m、
      フットプリント重なり \(\leq 10^{-9}\) \\
    共通尺度 & 元尺度の相対不一致 \(\leq 0.07290401\) \\
    対応構築 & 距離 \(\leq 0.25\) m、カメラあたり 3--200 対応 \\
    fit 受理 & カメラ \(\geq 6\)、対応 \(\geq 100\)、
      0.30 m 以内の内点率 \(\geq 90\%\)、RMS と q95 \(\leq 0.30\) m \\
    holdout 受理 & カメラ \(\geq 3\)、対応 \(\geq 80\)、
      0.30 m 以内の内点率 \(\geq 90\%\)、RMS と q95 \(\leq 0.30\) m \\
    全テンプレート & 0.30 m 以内の内点率 \(\geq 60\%\)、
      q95 \(\leq 1.50\) m、全意味セグメントの内点率 \(\geq 50\%\) \\
    変換と投影 & 逆変換の絶対許容 \(10^{-6}\)、
      投影の絶対許容 \(10^{-5}\) px \\
    \bottomrule
  \end{tabularx}
\end{table}


\subsection{Controlled 3D camera trajectories}
\label{sec:trajectories}

\begin{figure}[t]
  \centering
  \resizebox{\linewidth}{!}{% Controlled camera synthesis from path families to leakage-free partitions.
\begin{tikzpicture}[
  font=\sffamily\scriptsize,
  panel/.style={draw=black!28, rounded corners=2.5pt, fill=white,
    minimum height=50mm, inner sep=0pt, line width=0.45pt},
  title/.style={font=\sffamily\bfseries\footnotesize, text=black!84},
  tinylabel/.style={font=\sffamily\fontsize{6.35}{7.15}\selectfont,
    text=black!70, align=center},
  arr/.style={-{Stealth[length=1.8mm,width=1.25mm]}, line width=0.65pt},
  camera/.style={regular polygon, regular polygon sides=3, minimum size=3.4mm,
    inner sep=0pt, draw=gsblue, fill=gsblue!17, line width=0.5pt},
  sample/.style={circle, fill=gsblue, draw=white, line width=0.3pt,
    minimum size=1.9mm, inner sep=0pt},
  splitbox/.style={rounded corners=1.5pt, minimum width=31mm,
    minimum height=6.6mm, align=left, inner xsep=3pt, line width=0.5pt},
]
  % Four compact panels correspond to generation decisions made once per group.
  \node[panel, minimum width=43mm] (pa) at (2.15,0) {};
  \node[panel, minimum width=40mm] (pb) at (6.43,0) {};
  \node[panel, minimum width=42mm] (pc) at (10.65,0) {};
  \node[panel, minimum width=39mm] (pd) at (14.82,0) {};
  \node[title, anchor=north] at ($(pa.north)+(0,-2.2mm)$) {(a) Path family};
  \node[title, anchor=north] at ($(pb.north)+(0,-2.2mm)$) {(b) 3-D arc length};
  \node[title, anchor=north] at ($(pc.north)+(0,-2.2mm)$) {(c) Target + look-at};
  \node[title, anchor=north] at ($(pd.north)+(0,-2.2mm)$) {(d) Group-disjoint split};

  % (a) Horizontal footprints and vertical profiles are independently controlled.
  \node[tinylabel, anchor=west] at (0.35,1.52) {ground footprint};
  \draw[courtgreen, line width=0.75pt] (1.12,0.78) circle[radius=0.54];
  \fill[courtgreen] (1.12,0.78) circle[radius=0.75pt];
  \draw[gsblue, line width=0.75pt, rotate around={22:(3.08,0.78)}]
    (3.08,0.78) ellipse[x radius=0.78,y radius=0.42];
  \fill[gsblue] (3.08,0.78) circle[radius=0.75pt];
  \node[tinylabel, text=courtgreen] at (1.12,0.04) {circle};
  \node[tinylabel, text=gsblue] at (3.08,0.04) {ellipse};
  \draw[black!22] (0.36,-1.36) -- (4.00,-1.36);
  \node[tinylabel, anchor=west] at (0.72,-0.30) {camera height};
  \draw[courtgreen, line width=0.75pt] (0.58,-0.88) -- (1.88,-0.88);
  \draw[gsblue, line width=0.75pt]
    plot[smooth] coordinates {(2.22,-0.88) (2.48,-0.57) (2.75,-0.46)
      (3.03,-0.68) (3.30,-1.00) (3.58,-1.10) (3.88,-0.88)};
  \draw[arr, draw=black!48] (0.46,-1.31) -- (0.46,-0.55)
    node[above, tinylabel] {$z$};
  \node[tinylabel, text=courtgreen] at (1.23,-1.16) {planar};
  \node[tinylabel, text=gsblue] at (3.05,-1.39) {sinusoidal};
  \node[tinylabel, text width=37mm, anchor=south] at (2.15,-2.30)
    {$p(\theta)=[a\cos\theta,\,b\sin\theta,\,
      h+A\sin(k\theta+\phi)]^\top$};

  % (b) Samples are equally spaced along the full spatial curve, not in theta.
  \coordinate (o3) at (4.86,-1.18);
  \draw[arr, draw=black!45] (o3) -- ++(0.52,0) node[right, tinylabel] {$x$};
  \draw[arr, draw=black!45] (o3) -- ++(-0.27,0.34) node[above, tinylabel] {$y$};
  \draw[arr, draw=black!45] (o3) -- ++(0,0.60) node[above, tinylabel] {$z$};
  \draw[gsblue, line width=0.9pt]
    plot[smooth cycle, tension=0.55] coordinates {
      (5.10,-0.46) (5.31,0.16) (5.91,0.72) (6.73,0.98)
      (7.47,0.72) (7.86,0.14) (7.69,-0.49) (7.05,-0.89)
      (6.28,-0.96) (5.55,-0.80)};
  \foreach \x/\y in {5.10/-.46,5.31/.16,5.91/.72,6.73/.98,
      7.47/.72,7.86/.14,7.69/-.49,7.05/-.89,6.28/-.96,5.55/-.80}
    \node[sample] at (\x,\y) {};
  \draw[black!30, dashed] (5.91,0.72) -- (5.91,-1.22)
    (7.47,0.72) -- (7.47,-1.22);
  \draw[<->, >=Stealth, black!55, line width=0.5pt]
    (5.95,0.86) -- node[above, tinylabel] {$\Delta s$} (6.67,1.08);
  \node[tinylabel, fill=white, inner sep=1pt] at (6.43,-1.28)
    {equal spatial steps};
  \node[tinylabel, text width=35mm, anchor=south] at (6.43,-2.30)
    {$s(\theta)=\int_0^\theta\!\|p'(u)\|_2\,du$,\quad
     $s_i=iL/N$\\dense curve $\rightarrow$ invert cumulative $s$};

  % (c) A complex-centred sample binds to the metrically nearest court.
  \draw[courtgreen, fill=courtgreen!6, line width=0.48pt]
    (8.85,-0.82) rectangle (10.13,0.12);
  \draw[courtgreen, line width=0.42pt] (9.49,-0.82) -- (9.49,0.12);
  \draw[courtgreen, fill=courtgreen!13, line width=0.65pt]
    (11.18,-0.82) rectangle (12.46,0.12);
  \draw[courtgreen, line width=0.42pt] (11.82,-0.82) -- (11.82,0.12);
  \fill[courtgreen] (9.49,-0.35) circle[radius=0.8pt]
    (11.82,-0.35) circle[radius=0.8pt];
  \node[tinylabel] at (9.49,-1.06) {$\Court_A$};
  \node[tinylabel, text=courtgreen] at (11.82,-1.06) {$\Court_B=c_i^*$};
  \node[camera, shape border rotate=-75] (vc) at (11.35,1.13) {};
  \node[camera, shape border rotate=-42, opacity=0.55] (vp) at (10.60,0.86) {};
  \node[camera, shape border rotate=-122, opacity=0.55] (vn) at (12.24,0.78) {};
  \draw[gsblue!45, line width=0.55pt] (vp) to[bend left=10] (vc)
    to[bend left=10] (vn);
  \draw[black!42, dashed] (vc) -- node[left, tinylabel] {$d_A$} (9.49,-0.35);
  \draw[arr, draw=poseorange] (vc) -- node[right, tinylabel, text=poseorange]
    {$d_B$} (11.82,-0.35);
  \draw[arr, draw=gsblue!72] (vp) -- (11.82,-0.35);
  \draw[arr, draw=gsblue!72] (vn) -- (11.82,-0.35);
  \node[tinylabel, text width=37mm, anchor=south] at (10.65,-2.30)
    {$c_i^*=\operatorname*{lexmin}\arg\min_c
      \|v_i-o_c\|_2$\\camera $+z$ axes look at $o_{c_i^*}$};

  % (d) Seeded allocation moves complete trajectory groups as indivisible units.
  \node[tinylabel, draw=black!28, fill=softgray!72, rounded corners=1.5pt,
    minimum width=31mm, minimum height=8.0mm] (groups) at (14.82,1.08)
    {trajectory groups $G_1,\ldots,G_n$\\[-0.1em]seeded shuffle};
  \coordinate (forktop) at (13.10,0.70);
  \node[splitbox, draw=courtgreen, fill=courtgreen!9] (train) at (14.82,0.18)
    {\textcolor{courtgreen}{\textbf{Train 80\%}}\quad $G_2,G_4,\ldots$};
  \node[splitbox, draw=gsblue, fill=gsblue!8] (val) at (14.82,-0.58)
    {\textcolor{gsblue}{\textbf{Val. 10\%}}\quad $G_7$};
  \node[splitbox, draw=poseorange, fill=poseorange!8] (test) at (14.82,-1.34)
    {\textcolor{poseorange}{\textbf{Test 10\%}}\quad $G_9$};
  \draw[black!52, line width=0.55pt] (groups.south) -- ++(0,-0.12) --
    (forktop) -- (13.10,-1.34);
  \draw[arr, draw=courtgreen] (13.10,0.18) -- (train.west);
  \draw[arr, draw=gsblue] (13.10,-0.58) -- (val.west);
  \draw[arr, draw=poseorange] (13.10,-1.34) -- (test.west);
  \node[tinylabel, text width=33mm, anchor=south] at (14.82,-2.30)
    {all frames and views of one $G_i$ stay together};
\end{tikzpicture}
}
  \caption{\textbf{Deterministic view generation.}  Candidate circle and
  ellipse paths are scaled from captured-camera offsets and may be planar or
  sinusoidal in height.  Dense interpolation makes sampling uniform in 3D arc
  length.  Each camera looks at a deterministically resolved target court, and
  whole trajectory groups---not individual frames---are assigned to the
  80/10/10 splits.}
  \label{fig:trajectories}
\end{figure}

The generator derives one orbit centre for the tennis complex and one for each
accepted court.  Its base radius is the 90th percentile of captured cameras'
horizontal offsets from that centre, making the proposal scale specific to the
reconstructed venue.  A typed Cartesian product then enumerates circle/ellipse
shape, radius scale, axis ratio, in-plane orientation, base height, and vertical
profile.  A local path has the form
\begin{align}
  \bm{\gamma}(\theta)_{xy}
    &=\mathbf{R}(\phi)
      \begin{bmatrix}a\cos\theta\\b\sin\theta\end{bmatrix},
      &a&=r_{90}\alpha, &b&=a\beta, \\
  \gamma_z(\theta)
    &=h_0+A\sin(n\theta+\varphi),
  \label{eq:trajectory}
\end{align}
where \(\beta=1\) gives a circle and \(A=0\) gives a planar path.  All heights
must remain positive.

Uniform increments of \(\theta\) are not uniform in distance on an ellipse or
on a vertically modulated path.  We therefore sample 4096 points, construct the
cumulative 3D arc-length function, interpolate uniformly in arc length, and
increase the sample count until the closed-loop adjacent step is no larger than
the configured maximum.  Candidate order and selection use a fixed seed and
explicit stable-field ordering.  A coverage objective retains diverse
coverage/semantic-support modes and trajectory groups under a proposal budget.

The singleton and per-camera target behavior below is activated explicitly with
\texttt{dataset/court=v2}; the backward-compatible pipeline selector remains
v1 (seven two-point channels and group-level target assignment).  This
generation-schema v2 is distinct from the Court Detection V3 consumer/pose
contract: v2 defines the serialized singleton labels, whereas V3 defines the
canonical target-court training authority derived after loading.  Under v2, a
court-centred path retains its centre court.  A complex-centred path chooses the nearest accepted court to
each camera centre in metric 3D; distance ties within \(10^{-9}\) m are broken
lexicographically.  Let \(\mathbf{c}\) be the camera centre and \(\mathbf{q}\)
the chosen court look-at point.  We form an OpenCV camera basis with
\begin{equation}
  \mathbf{f}=\frac{\mathbf{q}-\mathbf{c}}{\|\mathbf{q}-\mathbf{c}\|_2},\quad
  \mathbf{r}=\mathbf{f}\times\mathbf{u}_{\mathrm{world}},\quad
  \mathbf{d}=\mathbf{f}\times\mathbf{r},
  \label{eq:lookat}
\end{equation}
so the axes are right, down, and forward.  For horizontal field of view
\(\omega\) and output width \(W\),
\begin{equation}
  f_x=f_y=\frac{W}{2\tan(\omega/2)},\qquad
  (c_x,c_y)=\left(\frac{W-1}{2},\frac{H-1}{2}\right).
  \label{eq:intrinsics}
\end{equation}
The complete path is sampled before split assignment.  All frames sharing a
trajectory-group identifier stay in one train, validation, or test split,
preventing near-duplicate neighbouring views from leaking across splits.

\begin{table}[t]
  \centering
  \caption{\textbf{軌道とサンプリングの方針。}
  数値は設定上の下限または予算であり、
  実際に採用された視点分布の証明ではない。
  採用後の分布は manifest から集計して報告する。
  }
  \label{tab:data_policy}
  \small
  \begin{tabularx}{\linewidth}{@{}p{0.26\linewidth}X@{}}
    \toprule
    次元 & 設定値 \\
    \midrule
    平面経路 & 円・楕円・矩形・超楕円。軸比 \(1.0,0.8,0.65\)。
      方位 \(0^\circ,45^\circ,90^\circ\) \\
    半径と高さ & 復元オフセットの \(0.65\)--\(0.95\) 倍。
      基準高さ \(1.5,2.25,3.0\) m \\
    鉛直経路 & 平面または正弦。振幅 \(0.0,0.25,0.5\) m \\
    視野 & 注視点高さ \(0\) または \(1.5\) m。
      水平画角 \(45^\circ\)--\(110^\circ\)（会場別の設定） \\
    サンプリング & seed 695。三次元弧長一様。隣接間隔の上限 1.05 m \\
    被覆 & 提案予算 4800。軌道グループ下限 24、採用フレーム下限 2000 \\
    分割と描画 & 軌道グループ単位で 8:1:1。8 shard \\
    SfM 制約 & 凸包を 0.5 m 内側へオフセット。空間被覆セル 1.0 m \\
    \bottomrule
  \end{tabularx}
\end{table}


\subsection{Rendering and exact labels}
\label{sec:labels}

The metric camera is converted to normalized scene units and passed to the
public 3DGS renderer.  The renderer returns RGB, depth, alpha, and pixel-support
evidence;
the court pipeline does not paste a separately rendered court onto the image.
Consequently the background, surface appearance, and occlusion structure
remain those reconstructed from the original video.

For canonical physical point \(\mathbf{X}_k\), the generator computes
\begin{equation}
  \mathbf{X}^{\Cam}_k =
  \T{\Cam}{\Scene_m}\T{\Scene_m}{\Court_i}\mathbf{X}_k,
  \qquad
  \tilde{\mathbf{u}}_k=\Kmat\mathbf{X}^{\Cam}_k,
  \qquad
  \mathbf{u}_k=\left(\frac{\tilde u_k}{\tilde w_k},
                       \frac{\tilde v_k}{\tilde w_k}\right).
  \label{eq:label_projection}
\end{equation}
Each record retains metric scene position, camera depth, pixel/normalized
coordinates, in-front state, in-frame state, and a renderer-support flag.  That
flag means that a radius-one pixel neighbourhood contains sufficient rendered
alpha and positive rendered depth; it is not a metric point-depth occlusion
test.  The
consumer contract reorders the fourteen physical points into camera-relative
near/far singleton channels but never guesses identity from image coordinates.
Segmentation and line targets are materialized from the same projected court
geometry; training-time crop, affine/perspective, color, blur, geometric
validity, and supervision-eligibility logic update all supervised quantities in
one augmented image frame.

The published sample stores \(\T{\Scene_m}{\Cam}\),
\(\T{\Scene_m}{\Court_i}\), and \(\Kmat\) rather than duplicating a pose10D
field.  After image augmentation, the model adapter derives the camera label
\begin{equation}
  \T{\Court_i}{\Cam}=
  \bigl(\T{\Scene_m}{\Court_i}\bigr)^{-1}\T{\Scene_m}{\Cam}.
  \label{eq:camera_to_court}
\end{equation}
It contains three translational and three rotational degrees of freedom.  The
adapter then applies the target-side canonicalization in
\cref{eq:canonical_pose} and constructs the pose10D target; augmented focal
length is carried separately from the rigid transform because the network
jointly predicts it for reprojection.  This is the paper's meaning of ``6DoF
supervision''; it must not be confused with the yaw-only player-location label
used by another task.

Once \cref{eq:camera_to_court} is accepted for a scene, another rendered view
requires no clicked points, line tracing, PnP solve, or per-frame extrinsic
calibration.  Human effort is concentrated in acquiring a suitable video and
auditing the scene-level alignment.  This amortization is what makes the
global-pose objective in \cref{sec:model,sec:consistency} practical at scale.
